\documentclass[12pt,CJK]{article}
\usepackage{geometry}
\usepackage{CJKutf8}
\usepackage{amsmath}
\usepackage{amssymb}
\usepackage{graphicx}
\usepackage{eufrak}
\usepackage{color}
\usepackage{slashed}
\usepackage{simplewick}
\usepackage{tikz}
\usepackage{tcolorbox}
\graphicspath{{../figures/}}
%%figures
\def\lfig#1#2{\includegraphics[width=#1 in]{#2}}
\def\addfig#1#2{\begin{center}\includegraphics[width=#1 in]{#2}\end{center}}
\def\wulian{\includegraphics[width=0.18in]{emoji_wulian.jpg}}
\def\bigwulian{\includegraphics[width=0.35in]{emoji_wulian.jpg}}
\def\bye{\includegraphics[width=0.18in]{emoji_bye.jpg}}
\def\bigbye{\includegraphics[width=0.35in]{emoji_bye.jpg}}
\def\huaixiao{\includegraphics[width=0.18in]{emoji_huaixiao.jpg}}
\def\bighuaixiao{\includegraphics[width=0.35in]{emoji_huaixiao.jpg}}
\def\jianxiao{\includegraphics[width=0.18in]{emoji_jianxiao.jpg}}
\def\bigjianxiao{\includegraphics[width=0.35in]{emoji_jianxiao.jpg}}
%% colors
\def\blacktext#1{{\color{black}#1}}
\def\bluetext#1{{\color{blue}#1}}
\def\redtext#1{{\color{red}#1}}
\def\darkbluetext#1{{\color[rgb]{0,0.2,0.6}#1}}
\def\skybluetext#1{{\color[rgb]{0.2,0.7,1.}#1}}
\def\cyantext#1{{\color[rgb]{0.,0.5,0.5}#1}}
\def\greentext#1{{\color[rgb]{0,0.7,0.1}#1}}
\def\darkgray{\color[rgb]{0.2,0.2,0.2}}
\def\lightgray{\color[rgb]{0.6,0.6,0.6}}
\def\gray{\color[rgb]{0.4,0.4,0.4}}
\def\blue{\color{blue}}
\def\red{\color{red}}
\def\green{\color{green}}
\def\darkgreen{\color[rgb]{0,0.4,0.1}}
\def\darkblue{\color[rgb]{0,0.2,0.6}}
\def\skyblue{\color[rgb]{0.2,0.7,1.}}
%%control
\def\be{\begin{equation}}
\def\ee{\nonumber\end{equation}}
\def\bea{\begin{eqnarray}}
\def\eea{\nonumber\end{eqnarray}}
\def\bch{\begin{CJK}{UTF8}{gbsn}}
\def\ech{\end{CJK}}
\def\bitem{\begin{itemize}}
\def\eitem{\end{itemize}}
\def\bcenter{\begin{center}}
\def\ecenter{\end{center}}
\def\bex{\begin{minipage}{0.2\textwidth}\includegraphics[width=0.6in]{jugelizi.png}\end{minipage}\begin{minipage}{0.76\textwidth}}
\def\eex{\end{minipage}}
\def\chtitle#1{\frametitle{\bch#1\ech}}
\def\bmat#1{\left(\begin{array}{#1}}
\def\emat{\end{array}\right)}
\def\bcase#1{\left\{\begin{array}{#1}}
\def\ecase{\end{array}\right.}
\def\bmini#1{\begin{minipage}{#1\textwidth}}
\def\emini{\end{minipage}}
\def\tbox#1{\begin{tcolorbox}#1\end{tcolorbox}}
\def\pfrac#1#2#3{\left(\frac{\partial #1}{\partial #2}\right)_{#3}}
%%symbols
\def\bropt{\,(\ \ \ )}
\def\optlist#1#2#3{(A)\;#1\ \ \ (B)\;#2\ \ \ (C)\;#3}
\def\sone{$\star$}
\def\stwo{$\star\star$}
\def\sthree{$\star\star\star$}
\def\sfour{$\star\star\star\star$}
\def\sfive{$\star\star\star\star\star$}
\def\rint{{\int_\leftrightarrow}}
\def\roint{{\oint_\leftrightarrow}}
\def\stdHf{{\textit{\r H}_f}}
\def\deltaH{{\Delta \textit{\r H}}}
\def\ii{{\dot{\imath}}}
\def\skipline{{\vskip0.1in}}
\def\skiplines{{\vskip0.2in}}
\def\lagr{{\mathcal{L}}}
\def\hamil{{\mathcal{H}}}
\def\vecv{{\mathbf{v}}}
\def\vecx{{\mathbf{x}}}
\def\vecy{{\mathbf{y}}}
\def\veck{{\mathbf{k}}}
\def\vecp{{\mathbf{p}}}
\def\vecn{{\mathbf{n}}}
\def\vecA{{\mathbf{A}}}
\def\vecP{{\mathbf{P}}}
\def\vecsigma{{\mathbf{\sigma}}}
\def\hatJn{{\hat{J_\vecn}}}
\def\hatJx{{\hat{J_x}}}
\def\hatJy{{\hat{J_y}}}
\def\hatJz{{\hat{J_z}}}
\def\hatj#1{\hat{J_{#1}}}
\def\hatphi{{\hat{\phi}}}
\def\hatq{{\hat{q}}}
\def\hatpi{{\hat{\pi}}}
\def\vel{\upsilon}
\def\Dint{{\mathcal{D}}}
\def\adag{{\hat{a}^\dagger}}
\def\bdag{{\hat{b}^\dagger}}
\def\cdag{{\hat{c}^\dagger}}
\def\ddag{{\hat{d}^\dagger}}
\def\hata{{\hat{a}}}
\def\hatb{{\hat{b}}}
\def\hatc{{\hat{c}}}
\def\hatd{{\hat{d}}}
\def\hatN{{\hat{N}}}
\def\hatH{{\hat{H}}}
\def\hatp{{\hat{p}}}
\def\Fup{{F^{\mu\nu}}}
\def\Fdown{{F_{\mu\nu}}}
\def\newl{\nonumber \\}
\def\vece{\mathrm{e}}
\def\calM{{\mathcal{M}}}
\def\calT{{\mathcal{T}}}
\def\calR{{\mathcal{R}}}
\def\barpsi{\bar{\psi}}
\def\baru{\bar{u}}
\def\barv{\bar{\upsilon}}
\def\qeq{\stackrel{?}{=}}
\def\torder#1{\mathcal{T}\left(#1\right)}
\def\rorder#1{\mathcal{R}\left(#1\right)}
\def\contr#1#2{\contraction{}{#1}{}{#2}#1#2}
\def\trof#1{\mathrm{Tr}\left(#1\right)}
\def\trace{\mathrm{Tr}}
\def\comm#1{\ \ \ \left(\mathrm{used}\ #1\right)}
\def\tcomm#1{\ \ \ (\text{#1})}
\def\slp{\slashed{p}}
\def\slk{\slashed{k}}
\def\calp{{\mathfrak{p}}}
\def\veccalp{\mathbf{\mathfrak{p}}}
\def\Tthree{T_{\tiny \textcircled{3}}}
\def\pthree{p_{\tiny \textcircled{3}}}
\def\dbar{{\,\mathchar'26\mkern-12mu d}}
\def\erf{\mathrm{erf}}
\def\const{\mathrm{constant}}
\def\pheat{\pfrac p{\ln T}V}
\def\vheat{\pfrac V{\ln T}p}
%%units
\def\fdeg{{^\circ \mathrm{F}}}
\def\cdeg{^\circ \mathrm{C}}
\def\atm{\,\mathrm{atm}}
\def\angstrom{\,\text{\AA}}
\def\SIL{\,\mathrm{L}}
\def\SIkm{\,\mathrm{km}}
\def\SIyr{\,\mathrm{yr}}
\def\SIGyr{\,\mathrm{Gyr}}
\def\SIV{\,\mathrm{V}}
\def\SImV{\,\mathrm{mV}}
\def\SIeV{\,\mathrm{eV}}
\def\SIkeV{\,\mathrm{keV}}
\def\SIMeV{\,\mathrm{MeV}}
\def\SIGeV{\,\mathrm{GeV}}
\def\SIcal{\,\mathrm{cal}}
\def\SIkcal{\,\mathrm{kcal}}
\def\SImol{\,\mathrm{mol}}
\def\SIN{\,\mathrm{N}}
\def\SIHz{\,\mathrm{Hz}}
\def\SIm{\,\mathrm{m}}
\def\SIcm{\,\mathrm{cm}}
\def\SIfm{\,\mathrm{fm}}
\def\SImm{\,\mathrm{mm}}
\def\SInm{\,\mathrm{nm}}
\def\SImum{\,\mathrm{\mu m}}
\def\SIJ{\,\mathrm{J}}
\def\SIW{\,\mathrm{W}}
\def\SIkJ{\,\mathrm{kJ}}
\def\SIs{\,\mathrm{s}}
\def\SIkg{\,\mathrm{kg}}
\def\SIg{\,\mathrm{g}}
\def\SIK{\,\mathrm{K}}
\def\SImmHg{\,\mathrm{mmHg}}
\def\SIPa{\,\mathrm{Pa}}

\geometry{tmargin=0.3in, bmargin=0.5in, lmargin=0.5in, rmargin=0.9in, nohead, nofoot}
\def\mark#1{{\color{blue} (#1分)}}
\renewcommand{\thepage}{}
\begin{document}
\bch
{\large 数理方法 课堂小测III 初入江湖版}

{\vskip 0.2in}

特殊函数定义：
\bea
J_\nu(x) &=& \sum_{k=0}^\infty \frac{(-1)^k}{k!(k+\nu)!}\left(\frac{x}{2}\right)^{2k+\nu}; \newl
Y_\nu(x) &=& \lim_{\mu\rightarrow \nu}\frac{J_\mu(x)\cos\mu\pi - J_{-\mu}(x)}{\sin (\mu\pi)}; \newl
j_\ell(x) &=& \sqrt{\frac{\pi}{2x}} J_{\ell+1/2}(x),\ \ \ell\in Z; \newl
y_\ell(x) &=& \sqrt{\frac{\pi}{2x}} Y_{\ell+1/2}(x), \ \ \ell\in Z; \newl
P_{\ell}(x) &=& \sum_{k=0}^\ell \frac{(\ell + k)!}{(k!)^2(\ell-k)!}\left(\frac{x-1}{2}\right)^k,\ \ \ell\in Z, \ell\ge 0; \newl
Y_{\ell, m}(\theta,\phi) &=&\frac{1}{2^\ell \ell !}\sqrt{\frac{(2\ell+1)}{4\pi} \frac{(\ell-m)!}{(\ell+m)!}}\left[\sin^m\theta \left(\frac{1}{\sin\theta}\frac{d}{d\theta}\right)^{\ell+m}\sin^{2\ell}\theta \right] e^{\ii m\phi},\ \ \ell,m\in Z, \ell\ge|m|.
\eea

\bitem
\item[(一)]{选择题，每题3分，共30分。

  \bitem

\item[(1)]{下列哪个物理问题对应的方程{\bf 不是}波动方程 \bropt
  
  \foptlist{均匀弦的横向小振动}{均匀弹性杆的纵向小振动}{孤立均匀球的温度变化}{引力波的传播}}


\item[(2)]{下列哪个函数具有分离变量的形式 \bropt
  
  \optlist{$e^{-xt}$}{$\frac{x^2}{1+e^t}$}{$\sqrt{t^2+x^2}$}{$e^{\frac{x}{t}}$}}
  

\item[(3)]{下列关于格林函数的说法{\bf 错误}的是 \bropt
  
  \optlist{是线性系统对脉冲输入的响应}{依赖于脉冲输入的位置}{依赖于边界条件}{只适用于无限大空间的问题}}


\item[(4)]{球坐标系$(r,\theta,\phi)$的拉普拉斯算符为： \bropt
  
  \goptlist{$\frac{1}{r^2}\frac{\partial}{\partial r}\left(r^2\frac{\partial}{\partial r}\right) + \frac{1}{r^2\sin\theta}\frac{\partial}{\partial \theta}\left(\sin\theta\frac{\partial}{\partial\theta}\right) + \frac{1}{r^2\sin^2\theta}\frac{\partial^2}{\partial \phi^2} $}
           {$\frac{1}{r}\frac{\partial}{\partial r}\left(r\frac{\partial}{\partial r}\right) + \frac{1}{r^2}\frac{\partial^2}{\partial \theta^2} + \frac{1}{r^2\sin^2\theta}\frac{\partial^2}{\partial \phi^2} $}
           {$\frac{\partial^2}{\partial r^2} + \frac{\partial^2}{\partial \theta^2} +\frac{\partial^2}{\partial \phi^2}$}
           {$\frac{\partial^2}{\partial r^2} + \frac{1}{r^2}\frac{\partial^2}{\partial \theta^2} + \frac{1}{r^2\sin^2\theta}\frac{\partial^2}{\partial \phi^2}$}
}  

\item[(5)]{贝塞尔函数的{\bf 导函数} $J'_{100}(x)$ 的最小正实数零点大约为 \bropt

  \optlist{$1$}{$10$}{$100$}{$10^4$}}
  
\item[(6)]{ 下列哪个贝塞尔函数在 $x=0$ 处发散？ \bropt
  
  \optlist{$J_1(x)$}{$J_{-1}(x)$}{$J_{1/2}(x)$}{$J_{-1/2}(x)$} }

\item[(7)]{ 积分$\int_0^a J_1(x) dx= $ \bropt
  
  \optlist{$J_0(a)$}{$1-J_0(a)$}{$aJ_0(a)$}{$a J_{1}(a)$} }

\item[(8)]{两类贝塞尔函数的平方和 $[J_{5}(100)]^2+[Y_{5}(100)]^2$ 大约为 \bropt

  \optlist{$0.006$}{$0.01$}{$0.06$}{$1$}}  

\item[(9)]{设$\mu, \nu$是球贝塞尔函数$j_\ell(x)$ ($\ell\ge 0$)的两个不同的零点，则下列哪个等式一定正确？ \bropt
  
  \goptlist{$\int_0^1j_\ell(\mu x)j_\ell(\nu x)dx = 0 $}{ $\int_0^1j_\ell(\mu x)j_\ell(\nu x)xdx = 0 $}{$\int_0^1j_\ell(\mu x)j_\ell(\nu x)x^2dx = 0 $}{$\int_0^{1}(1-x^2)j_\ell(\mu x)j_\ell(\nu x)dx = 0 $}}
  
\item[(10)]{$\cos\theta\, Y_{4,2}(\theta,\phi)$ 可写成哪两个球面谐函数的线性组合？ \bropt
  
  \foptlist{$ Y_{3,2}(\theta,\phi)$ 和 $Y_{5,2}(\theta,\phi)$}{$Y_{4,2}(\theta,\phi)$ 和 $Y_{5,2}(\theta,\phi)$}{$Y_{3,2}(\theta,\phi)$ 和 $Y_{4,2}(\theta,\phi)$}{$Y_{4,1}(\theta,\phi)$ 和 $Y_{4,3}(\theta,\phi)$}}  

\eitem  
}
\item[(二)]{填空题（每题5分，共30分）
  \bitem
\item[(1)]{设二维正交曲面坐标系$(x,y)$的正交线元长度依次为 $dx$ 和 $e^{-x^2}dy$，写出该坐标系的拉普拉斯算符的显式表达: $\nabla^2 f= $\uline{3}。}
\item[(2)]{积分$\int_0^1 x^4 P_6(x) dx = $ \uline{1}。}
\item[(3)]{函数 $f_1(x,y)$, $f_2(x,y)$ 在椭圆 $2x^2+ y^2 = 1$ 上的取值处处为零，在该椭圆内部分别满足

  $ \nabla^2 f_1 =-k_1^2f_1$ 以及 $\nabla^2 f_2 = -k_2^2f_2$，其中$k_1>k_2>0$。

  则在该椭圆内的积分$\iint f_1(x,y)f_2(x,y) dx dy = $ \uline{1}。}
\item[(4)]{设 $\mu>0$，$J_2(\mu)=0$，则积分$\int_0^1\left[J_2(\mu x)\right]^2 x dx =$ \uline{1}。}  
\item[(5)]{一维无边界的空间内满足方程 $\frac{\partial^2u}{\partial t^2}-a^2\frac{\partial^2u}{\partial x^2} = 0$ 和初始条件  $\left.u\right\vert_{t=0} = e^{-x^2},\ \ \left.\frac{\partial u}{\partial t}\right\vert_{t=0} = 0$ 的解是: $u(x,t)=$ \uline{3}。 }
\item[(6)]{一维无边界的空间内满足方程 $\frac{\partial u}{\partial t}-a\frac{\partial^2u}{\partial x^2} = 0$ 和初始条件  $\left.u\right\vert_{t=0} = e^{-x^2}$ 的解是: $u(x,t)=$ \uline{3}。 }

  \eitem
}
\item[(三)]{两端固定的，长度为$L$的均匀弦的横向振动$u(x,t)$ ($0\le x\le L$)满足初始条件：
  \bea
  \left.u\right\vert_{t=0} &=& A\sin \frac{\pi x}{L} \newl
  \left. \frac{\partial u}{\partial t}\right\vert_{t=0} &=& 0.
  \eea
  求解$u(x,t)$。设弦的线密度$\lambda$，张力$f$均已知。（20分)
}

\item[(四)]{有半径为 $R$，导热系数为 $\lambda$，单位质量比热为 $c$，质量密度为 $\rho$ 的孤立均匀薄圆盘。以圆盘中心为原点建立极坐标 $(r,\theta)$。初始时刻各点的温度为 $\left.T\right\vert_{t=0}= T_0\left(2 + \frac{r^2}{R^2}\cos\theta\right)$。计算之后圆盘上各点的温度变化。
(20分)}

  
\eitem  





\ech
\end{document}
